\section[Cap\'{i}tulo 4: Amenazas Digitales Representables en el Videojuego]{\centering Cap\'{i}tulo 4: Amenazas Digitales Representables en el Videojuego}

Las amenazas digitales seleccionadas para el prototipo deben ser comprensibles para estudiantes y, al mismo tiempo, cercanas a situaciones reales. Por esa raz\'{o}n se priorizan casos donde el jugador pueda observar evidencias en pantalla, comparar datos y decidir una acci\'{o}n de respuesta (CISA, 2026).

Este cap\'{i}tulo ordena los tipos de amenazas que pueden convertirse en niveles o casos jugables. La selecci\'{o}n no busca cubrir todo el campo de ciberseguridad, sino trabajar incidentes frecuentes y adecuados para una primera versi\'{o}n del videojuego (OWASP Foundation, 2026).

\subsection{Amenazas base para escenarios}

\begin{center}
\begin{tabular}{|p{0.24\textwidth}|p{0.35\textwidth}|p{0.28\textwidth}|}
\hline
\textbf{Amenaza} & \textbf{Pistas visibles} & \textbf{Decisi\'{o}n esperada} \\
\hline
\textit{Phishing} & Dominio alterado, urgencia falsa, adjuntos dudosos o errores en el remitente (CISA, 2026). & Reportar, bloquear enlace o escalar el caso. \\
\hline
Malware & Archivo inesperado, extensi\'{o}n sospechosa o alerta del sistema (Stallings \& Brown, 2018). & Aislar equipo o rechazar ejecuci\'{o}n. \\
\hline
Ingenier\'{i}a social & Solicitud de datos, presi\'{o}n emocional o suplantaci\'{o}n de autoridad (Verizon, 2024). & Verificar identidad antes de actuar. \\
\hline
Ransomware & Bloqueo de archivos, nota de rescate o actividad an\'{o}mala (CISA, 2026). & Contener, informar y evitar propagaci\'{o}n. \\
\hline
\end{tabular}
\end{center}

\subsection{Relaci\'{o}n entre amenaza, pista y decisi\'{o}n}

Para que la experiencia funcione como aprendizaje pr\'{a}ctico, cada amenaza debe presentar pistas interpretables. Si la respuesta correcta aparece de forma demasiado evidente, el jugador solo sigue instrucciones; si no existen pistas suficientes, la decisi\'{o}n se vuelve azarosa (Prince, 2004).

El dise\~{n}o de casos debe equilibrar claridad y dificultad. Un correo puede incluir un dominio casi correcto, un archivo con nombre cre\'{i}ble y una justificaci\'{o}n urgente; el estudiante debe relacionar esos detalles antes de actuar, como ocurrir\'{i}a en una revisi\'{o}n real de incidentes (NIST, 2024).

La progresi\'{o}n del videojuego puede iniciar con se\~{n}ales directas y luego combinar varias amenazas dentro de un mismo caso. Esto permite que el jugador pase de reconocer patrones simples a analizar situaciones con informaci\'{o}n parcial o contradictoria (Ng \& Hasan, 2024).

\subsection{S\'{i}ntesis del cap\'{i}tulo}

Las amenazas seleccionadas permiten convertir conceptos de ciberseguridad en casos jugables. El valor del prototipo est\'{a} en presentar pistas, decisiones y consecuencias que obliguen al estudiante a razonar antes de responder.
